% Lösungen Einheit 2 von 5 – Eigenschaften aller Graphen am Term (zusammenbau v0.1)
\einheitenkopf{Einheit 2 von 5 · Eigenschaften aller Graphen am Term}

\erg{12}{a) fällt heraus \quad b) $x_1 = 0$, $x_2 = 6$; der Faktor $a$ ist nie null. \quad c) Für $x < 0$ ist $x^3 < 0$, und $e^{ax} > 0$ für jedes $a$; das Produkt ist negativ. \quad d) Der e-Faktor ist nie null; $x^2 = a$: für $a > 0$ zwei, für $a = 0$ eine, für $a < 0$ keine Nullstelle. \quad e) Für $a < 0$ gegen $0$, weil der e-Faktor überwiegt; für $a > 0$ und für $a = 0$ (dann $f_0(x) = 2x + 1$) gegen $+\infty$. \quad f) Der Parameterterm verschwindet für $x^2 - 5x = 0$, also $x = 0$ oder $x = 5$: gemeinsame Punkte $(0 | 3)$ und $(5 | 3)$. \quad g) $f_a(-x) = (-x)^3 \cdot e^{a(-x)^2} = -x^3 \cdot e^{a x^2} = -f_a(x)$ für jedes $a$. \quad h) $f_a'(x) = (2 + 2ax) \cdot e^{ax}$, also $f_a'(0) = 2$ für jedes $a$. \quad i) (1) Alle Graphen gehen durch $(1 | 2)$. (2) Dort haben alle dieselbe Steigung, also dieselbe Tangente. (3) Zwei verschiedene Graphen haben außer $(1 | 2)$ keinen gemeinsamen Punkt.}
\erg{13}{a) Richtig ist nur $x = a$: der Faktor $e^{x+1}$ ist nie null, $x = -1$ ist keine Nullstelle. \quad b) Ein gemeinsamer Punkt liegt auf allen Graphen, sein y-Wert darf also nicht von $a$ abhängen; aus $a_1 x^2 = a_2 x^2$ mit $a_1 \neq a_2$ folgt $x = 0$, und $f_a(0) = 0$ für jedes $a$.}
